\documentclass[10pt,a4paper]{amsart}
\usepackage[margin=19mm]{geometry}
\usepackage{amsmath,amssymb,mathtools}
\usepackage[scr=rsfs,cal=euler]{mathalfa}
\usepackage{xfrac}
\usepackage[unicode,hidelinks,bookmarks=false]{hyperref}
\newcommand{\ie}{i.e.\ }
\newcommand{\cf}{cf.\ }
\newcommand{\resp}{respectively\ }
\newcommand{\Subsection}{\subsection}
\providecommand{\nobreakdash}{\nobreak\hbox{-}}
\setlength{\emergencystretch}{2em}
\multlinegap0pt
\usepackage{bm}
\usepackage{bbm}
\usepackage{dsfont}
\usepackage{xspace}
\usepackage{cancel}
\usepackage{mathtools}
\usepackage{amsthm}
\makeatletter
\@ifundefined{theorem}{\newtheorem{theorem}{Theorem}}{}
\@ifundefined{lemma}{\newtheorem{lemma}[theorem]{Lemma}}{}
\makeatother
\title[Stability and blow-up for a suspension bridge plate model wi]{Stability and blow-up for a suspension bridge plate model with fractional damping and memory}
\author{Iqra Kanwal, Jianghao Hao, Muhammad Fahim Aslam, Zayd Hajjej, Mauricio Sep\'ulveda-Cort\'es, Rodrigo Vej\'ar-Asem}
\date{Source: arXiv:2603.20566v1 (CC BY 4.0)}
\begin{document}
\maketitle

\noindent\emph{Source and licence.} Stability and blow-up for a suspension bridge plate model with fractional damping and memory. Version arXiv:2603.20566v1 (CC BY 4.0), distributed under the Creative Commons Attribution 4.0 International licence, as recorded in the arXiv licence metadata for this exact version and shown on the abstract page https://arxiv.org/abs/2603.20566v1. Full rights notice: licenses/arxiv-2603.20566-CC-BY-4.0.txt.\par\smallskip

\noindent\textit{Editorial scope.} Counted unit: the Theorem whose source label is \texttt{teo6.2} in the article (source0.tex lines 613--619, source label \texttt{teo6.2}), delivered together with the article's own complete original proof of it (source0.tex lines 620--786, one \texttt{proof} environment of 167 lines). No mathematics is written, repaired or re-derived: statement and proof are the article's own text line for line. Required context retained uncounted: eq2.5 (lines 275--288); eq2.7 (lines 352--356); psi11 (lines 459--465); le5.1 (lines 605--610). Every definition, display and label of those lines is retained unchanged. Omitted from this package and not redistributed: 1--274, 289--351, 357--458, 466--604, 611--612, 787--1583. Nothing inside a retained excerpt is omitted or altered: every retained body is delivered exactly as the article's lines are written, so no omission receipt of any kind accompanies this package. Citations: the retained excerpts carry no citation command at all, so no bibliography entry of the article is transcribed and no reference is redistributed. Reference adaptation: none was needed and none was made; every reference inside the delivered unit resolves inside it, so no unresolved reference or citation is produced. Printed slips kept literal: no printed word, symbol, hypothesis or number of the counted statement or of its proof was altered. Layout and class: the article's own class is \texttt{article} with packages including inputenc, amssymb, amsfonts, amsmath, amsthm, mathrsfs, mathtools, color, hyperref, soul, xcolor; the wrapper is the lane's pinned \texttt{amsart} editorial wrapper at 10pt on A4, which restates the theorem-like environments the retained text needs and adds the article's own macro definitions that the retained text needs (none), with extra wrapper packages (bm, bbm, dsfont, xspace, cancel, mathtools). The delivered numbering is the unit's own. Assets: no retained span contains a figure environment, a graphics-inclusion command, a PDF-page inclusion, a table or a TikZ picture; no image file is copied into this package, the package declares no asset dependency and no image file is redistributed with it. Licence: version arXiv:2603.20566v1 of this article is distributed under the Creative Commons Attribution 4.0 International licence, as recorded in the arXiv licence metadata for this exact version and shown on the abstract page https://arxiv.org/abs/2603.20566v1. A full rights notice is delivered at licenses/arxiv-2603.20566-CC-BY-4.0.txt. Characters: every retained excerpt is pure ASCII, so the delivered engine, XeLaTeX with its default Latin Modern text font, reports no missing character.
\begin{equation} \label{eq2.5}
\hspace{-0.4cm}\left\{\begin{array}{lr}
u_{t t}(x, y, t)+\lambda\Delta^{2} u(x, y, t)+ \displaystyle\int_{0}^{+\infty} g(s) \Delta^2 \mu (x,y,t,s) ds \\[1em]
+\kappa \displaystyle\int_{-\infty}^{+\infty} \phi(x, y, \vartheta, t) \xi(\vartheta) d \vartheta=u|u|^{p-2},  & (x, y) \in \Omega, t>0, \\[1em]
 \phi_{t}(x, y, \vartheta, t)+\left(\vartheta^{2}+\beta\right) \phi(x, y, \vartheta, t)-u_t (x,y,t) \xi(\vartheta)=0, & (x, y) \in \Omega, \vartheta \in \mathbb{R}, t>0, \\[1em]
\mu_t (x,y,t,s) + \mu_s(x,y,t,s)=u_t(x,y,t),& (x, y) \in \Omega,\; s, t>0, \\[1em]
u(0, y, t)=u_{x x}(0, y, t)=u(\pi, y, t)=u_{x x}(\pi, y, t)=0, & (x, t) \in(0, \pi) \times(0,+\infty), \\[1em]
u_{y y}(x, \pm d, t)+\sigma u_{x x}(x, \pm d, t)=0, & (x, t) \in(0, \pi) \times(0,+\infty), \\[1em]
u_{y y y}(x, \pm d, t)+(2-\sigma) u_{x x y}(x, \pm d, t)=0, & (x, y) \in \Omega, t>0, \\[1em]
u(x, y, 0)=u_{0}(x, y), u_{t}(x, y, 0)=u_{1}(x, y), & (x, y) \in \Omega, \\[1em]
\phi(x, y, \vartheta, 0)=0, & (x, y) \in \Omega, \vartheta \in \mathbb{R}, \\[1em]
\mu(x,y,0) = 0, \quad \mu_0 (x,y,0,s) = u_0 (x) - u_0 (x,-s),   & (x, y) \in \Omega, s>0,
\end{array}\right.
\end{equation}

\begin{equation}\label{eq2.7}
E^{\prime}(t)
= \frac{1}{2} \int_{0}^{+\infty} g'(s) ||\mu(s)||^2_{H^2_{*}(\Omega))} ds
 -\kappa \int_{\Omega} \int_{-\infty}^{+\infty} (\vartheta^2 + \beta) |\phi(x,y,\vartheta,t)|^2 d\vartheta dx dy \leq 0.
\end{equation}

\begin{align}\label{psi11}
&-\int_{0}^{+\infty} g(s)\int_{\Omega} \Delta^{2}\mu(s) u(x,y,t) dx dyds\nonumber\\
&= -\int_{0}^{+\infty} g(s) ( \mu(s), u)_{H_*^{2}(\Omega)}\,ds \nonumber\\
&\le \frac{\lambda}{2}\|u\|^2_{H^2_*(\Omega)}
  + \frac{1-\lambda}{2\lambda}
    \int_{0}^{+\infty} g(s)\,\|\mu(s)\|_{H_*^{2}(\Omega)}^{2}\,ds.
\end{align}

\begin{lemma}\label{le5.1}
There exists $C_1>1$ such that
\begin{equation}
\|u\|_{p}^{s} \leq C_1\left(\|u\|_{p}^{p}+\|u\|^2_{H^2_*(\Omega)}\right)
\end{equation}
\end{lemma}

\begin{theorem}\label{teo6.2}
Assume that (H1)-(H2), $E(0)< 0$ and
\begin{equation}\label{eq5.3}
\displaystyle\int_{0}^{\infty} g(s) d s<\frac{p-2}{p},
\end{equation}
then the solution of system \eqref{eq2.5} blows-up in a finite time.
\end{theorem}

\begin{proof}
Let
\begin{equation*}
H(t)= -E(t).
\end{equation*}
By \eqref{eq2.7}, we have
\begin{eqnarray}\label{eq49}
H^{\prime}(t)&=&-E^{\prime}(t)\nonumber\\
&=& -\frac{1}{2} \int_0^{+\infty} g'(s) ||\mu(s)||^2_{H^2_{*}(\Omega)} ds\nonumber\\
&&\hspace{3cm}+\kappa \int_\Omega \int_{-\infty}^{+\infty} (\vartheta^2 + \beta) |\phi(x,y,\vartheta,t)|^2 d\vartheta dx dy \geq 0.
\end{eqnarray}
Thus, we get
\begin{equation}\label{eq5.6}
0 < H(0) \leq H(t) \leq \frac {1}{p}\|u\|_{p}^{p}.
\end{equation}
Let
\begin{equation} \label{eq51}
B(t) = H^{1-\gamma} (t) + \varepsilon \int_\Omega u u_t dx dy,
\end{equation}
where $\varepsilon > 0$ will be specified afterwards, and $\gamma$ is a positive constant that satisfies
\begin{equation}
0 < \gamma < \frac{p-2}{2p}.
\end{equation}
Differentiating (\ref{eq51}) and using \eqref{eq2.5}, we get
\begin{equation}
\begin{split}
B^{\prime}(t)
& =  \varepsilon  ||u_t||_2^2 - \varepsilon \lambda ||u||^2_{H^2_{*}(\Omega)}+\varepsilon\|u\|_{p}^{p}
\\
&\quad +(1 - \gamma) H^{-\gamma}(t) H'(t) -\kappa \varepsilon \int_\Omega u \int_{-\infty}^{+\infty}  \phi(x,y,\vartheta,t)\xi(\vartheta) d\vartheta dx dy
\\
&\quad -\varepsilon \displaystyle\int_{0}^{\infty} g(s) (u, \mu(s))_{H^2_{*}(\Omega)} ds.
\end{split}
\label{eq53}
\end{equation}
As in \eqref{psi11}, we infer that
\begin{equation}\label{eq55}
\int_{0}^{+\infty} g(s) ( \mu(s), u)_{H_*^{2}(\Omega)}ds\le (1-\lambda)\|u\|^2_{H^2_*(\Omega)}
  + \frac{1}{4}\int_{0}^{+\infty} g(s)\,\|\mu(s)\|_{H_*^{2}(\Omega)}^{2}\,ds.
\end{equation}
Substituting (\ref{eq55}) in (\ref{eq53}), we get
\begin{equation}
\begin{split}
B^{\prime}(t)
& \geq (1-\gamma) H^{-\gamma} (t)  H' (t)+ \varepsilon ||u_t||_2^2 -\varepsilon ||u||^2_{H_*^{2}(\Omega)}+\varepsilon\|u\|_{p}^{p}
\\
&\quad -\kappa\varepsilon  \int_\Omega u \int_{-\infty}^{+\infty}  \phi(x,y,\vartheta,t)\xi (\vartheta) d\vartheta dx dy
\\
&\quad -\frac{\varepsilon}{4} \int_{0}^{+\infty} g(s) ||\mu(s)||^2_{H^2_{*}(\Omega)} ds.
\end{split}
\label{eq56}
\end{equation}
By the help of (\ref{eq49}) and Young's inequality, we have, for any $\delta > 0$,
\begin{equation}
\begin{split}
\kappa
& \int_\Omega u \int_{-\infty}^{+\infty}  \phi(x,y,\vartheta,t)\xi (\vartheta) d\vartheta dx dy
\\
&\quad \leq \delta \kappa M ||u||_2^2 + \frac{\kappa}{4\delta}  \int_\Omega \int_{-\infty}^{+\infty} (\vartheta^2 + \beta) |\phi (x,y,\vartheta,t)|^2 d\vartheta dx dy
\\
&\quad \leq \delta \kappa M ||u||_2^2 + \frac{1}{4\delta} H'(t),
\end{split}
\label{eq57}
\end{equation}
Inserting (\ref{eq57}) in (\ref{eq56}), we have
\begin{equation}
\begin{split}
B^{\prime}(t)
& \geq \biggl( (1-\gamma) H^{-\gamma}(t) - \frac{\varepsilon}{4\delta} \biggl)  H' (t)
\\
&\quad + \varepsilon ||u_t||_2^2 - \varepsilon ||u||^2_{H^2_{*}(\Omega)} - \varepsilon \delta \kappa M ||u||_2^2
\\
&\quad - \frac{\varepsilon}{4} \int_{0}^{+\infty} g(s) ||\mu(s)||^2_{H^2_{*}(\Omega)} ds+\varepsilon\|u\|_{p}^{p}.
\end{split}
\label{eq58}
\end{equation}
Next, by selecting $\delta$ such that
\begin{equation}
    \frac{1}{4\delta} = k H^{-\gamma}(t),
    \label{eq59}
\end{equation}
where $k$ is a positive constant to be determined later, (\ref{eq58}) becomes
\begin{equation}
\begin{split}
B^{\prime}(t)
& \geq \bigr[ (1-\gamma) - \varepsilon k \bigr] H^{-\gamma}(t)H'(t) - \varepsilon ||u||^2_{H^2_{*}(\Omega)} - \frac{\varepsilon \kappa M}{4k} H^{\gamma}(t)||u||_2^2\\
&\quad + \varepsilon ||u_t||_2^2+\varepsilon\|u\|_{p}^{p} - \frac{\varepsilon}{4} \int_{0}^{+\infty} g(s) ||\mu(s)||^2_{H^2_{*}(\Omega)} ds.
\label{eq60}
\end{split}
\end{equation}
Using \eqref{eq5.6}, we have
\begin{equation*}
H^{\gamma}(t) \leq \frac{1}{p^{\gamma}}\|u\|_{p}^{p \gamma},
\end{equation*}
which implies that
\begin{equation}\label{eq6Z}
\kappa M H^{\gamma}(t)\|u\|_{2}^{2} \leq C_{2}\|u\|_{p}^{p \gamma+2}
\end{equation}
for some $C_{2}>0$. Putting \eqref{eq6Z} in \eqref{eq60} and using Lemma \ref{le5.1} with $2\leq s=p \gamma+2 \leq p$, we obtain, for any $0<m<1$,
\begin{equation}
\begin{aligned}
B^{\prime}(t) \geq & ((1-\gamma)-\varepsilon k) H^{-\gamma}(t) H^{\prime}(t)+\varepsilon\frac{p(1-m)+2}{2}\left\|u_{t}\right\|_{2}^{2} \\
& +\varepsilon\left(m-\frac{C_{3}}{4 k}\right)\|u\|_{p}^{p}+\varepsilon\left(\frac{p(1-m)\lambda -2}{2}-\frac{C_{3}}{4k}\right)\|u\|^{2}_{H^2_{*}(\Omega)} \\
& +\varepsilon\frac{\kappa p(1-m)}{2}  \int_{\Omega} \int_{-\infty}^{+\infty}|\phi(x,y,\vartheta,t)|^{2} d\vartheta dx dy+\varepsilon p(1-m)H(t)\\
& +\varepsilon\frac{2p(1-m)-1}{4}\int_{0}^{+\infty} g(s)\left\|\mu(s)\right\|^{2}_{H^2_{*}(\Omega)} ds
\end{aligned}
\end{equation}
where $C_{3}=C_1 C_{2}$. \\
First, we choose $0<m<1$ such that
$$m<\min\Big\{1-\frac{1}{2p}, 1-\frac{2}{p\lambda} \Big\}.$$
This choice is justified by the assumption $p>2$ together with \eqref{eq5.3}.

Next, we select $k$ sufficiently large such that
$$m-\frac{C_{3}}{4 k}>0,\;\text{and}\;\; \frac{p(1-m)\lambda -2}{2}-\frac{C_{3}}{4k}>0.$$
Finally, we pick up $\varepsilon$ small enough so that
$$(1-\gamma) - \varepsilon k > 0,\;\;\text{and}\;\;H^{1-\gamma}(0)+\varepsilon \int_{\Omega} u_{0} u_{1} dx dy>0.
$$
Consequently, there exists a positive constant $C_4$ such that
\begin{align}\label{eq5.21}
B^{\prime}(t)&\geq C_4\left\{H(t)+\left\|u_{t}\right\|_{2}^{2}+\|u\|^{2}_{H^2_{*}(\Omega)}+\|u\|_{p}^{p}\right.\nonumber \\
&\quad\;\;\;\;\left.+\int_{\Omega} \int_{-\infty}^{+\infty}|\phi(x,y,\vartheta,t)|^{2} d\vartheta dx dy+\int_{0}^{+\infty} g(s)\left\|\mu(s)\right\|^2_{H^2_{*}(\Omega)} ds\right\}.
\end{align}
and
\begin{equation}\label{GrindEQ__66_}
B(t)\geq B(0)>0,\;\forall\;t>0.
\end{equation}
On the other hand, it is easy to see that
\begin{equation}\label{GrindEQ__67_}
\begin{aligned}
B^{\frac{1}{1-\gamma}}(t)= & \left[H^{1-\gamma}(t)+\varepsilon \int_{\Omega} u u_{t} dx dy \right]^{\frac{1}{1-\gamma}}
\end{aligned}
\end{equation}
\begin{equation}
    \leq C_5\Big[H(t)+\left|\int_{\Omega} u u_{t} dx dy\right|^{\frac{1}{1-\gamma}}\Big]
    \label{eq68}
\end{equation}
for some positive constant $C_5$.\\
Next, employing H\"{o}lder's inequality to obtain
\begin{equation*}\label{eq69}
\left| \int_{\Omega} u u_{t} dx dy \right| \leq\|u\|_{2} \left\|u_{t}\right\|_{2} \leq C_6 \|u\|_{p} \left\|u_{t}\right\|_{2}.
\end{equation*}
Therefore, by using Young's inequality and Lemma \ref{le5.1} with $2<s=\frac{2}{1-2\gamma}<p$, we have
\begin{eqnarray}\label{eq699}
\left|\int_{\Omega} u u_{t} dx dy\right|^{\frac{1}{1-\gamma}} &\leq& C^{\frac{1}{1-\gamma}}_6\|u\|_{p}^{\frac{1}{1-\gamma}} \left\|u_{t}\right\|_{2}^{\frac{1}{1-\gamma}} \nonumber\\
&\leq& C_7\left[\|u\|_{p}^{\frac{2}{1-2\gamma}}+\left\|u_{t}\right\|_{2}^2\right]\nonumber\\
&\leq& C_8 \left[ \|u\|^2_{H^2_*(\Omega)}+\|u\|_{p}^{p}+\left\|u_{t}\right\|_{2}^{2}\right].
\end{eqnarray}
where $C_6, C_7, C_8>0$. Putting \eqref{eq699} in (\ref{eq68}), we get the existence of a positive constant $C_9$ such that
\begin{equation}\label{eq72}
B^{\frac{1}{1-\gamma}}(t) \leq C_9\left[H(t)+\|u\|_{H^2_*(\Omega)}^{2}+\|u\|_{p}^{p}+\|u_{t}\|_{2}^{2}\right].
\end{equation}
From \eqref{eq5.21} and (\ref{eq72}), we deduce that
\begin{equation}
B^{\prime}(t) \geq \varpi B^{\frac{1}{1-\gamma}}(t),\;\forall\;t\geq 0,
\label{eq73}
\end{equation}
for some positive constant $\varpi$.
A simple integration of (\ref{eq73}) gives
$$
B^{\frac{\gamma}{1-\gamma}}(t) \geq \frac{1}{B^{\frac{-\gamma}{1-\gamma}}(0)-\varpi \frac{\gamma}{(1-\gamma)} t}.
$$
Consequently, $B(t)$ blows-up in finite time
$$
T \leq T^{*}=\frac{1-\gamma}{\varpi \gamma B^{\frac{\gamma}{1-\gamma}(0)}}.
$$
This completes the proof.
\end{proof}

\end{document}
